\documentclass[letterpaper,twocolumn,9pt]{extarticle}
\usepackage[letterpaper,top=0.62in,bottom=0.62in,left=0.66in,right=0.66in,columnsep=0.18in]{geometry}
\usepackage{times}
\usepackage{amsmath,amssymb,amsfonts}
\usepackage{graphicx}
\usepackage{booktabs}
\usepackage{url}
\usepackage{xcolor}
\usepackage{cite}
\usepackage[compact]{titlesec}
\usepackage{caption}
\captionsetup{labelsep=period,font=small,justification=centering,skip=3pt}
\captionsetup[figure]{name=Figure}
\captionsetup[table]{name=Table}
\graphicspath{{../08_figures/}}

\setlength{\parindent}{0.167in}
\setlength{\parskip}{0pt}
\setlength{\textfloatsep}{6pt plus 2pt minus 2pt}
\setlength{\floatsep}{6pt plus 2pt minus 2pt}
\setlength{\intextsep}{6pt plus 2pt minus 2pt}
\setlength{\abovecaptionskip}{3pt}
\setlength{\belowcaptionskip}{2pt}
\pagestyle{empty}
\renewcommand{\thesection}{\Roman{section}}
\renewcommand{\thesubsection}{\Alph{subsection}}
\titleformat{\section}[block]{\centering\fontsize{11}{13}\selectfont\scshape}{\thesection.}{0.5em}{}
\titlespacing*{\section}{0pt}{6pt}{4pt}
\titleformat{\subsection}[block]{\itshape\fontsize{9.5}{11}\selectfont}{\thesubsection.}{0.5em}{}
\titlespacing*{\subsection}{0pt}{5pt}{3pt}

% Metadata for converters (pandoc -> .docx). LaTeX ignores these: \maketitle is
% never called; the visual title block is the \twocolumn[...] group below.
\title{A License-Clean Graph Neural Network for Fast Parasitic Extraction in
PCB-Embedded Planar Magnetics}
\author{Duong Viet Hoang\thanks{Corresponding author} \and Lun-Min Shih\\
Department of Computer Science, Da-Yeh University, Taiwan\\
Hoangduong4316@icloud.com, Lmshih@mail.dyu.edu.tw}
\date{Version 6: v3 with the target naming corrected -- C_ps is the parasitic;
L_p,L_s,M are air-core inductances composed from FastHenry's partial-inductance
matrix. Supersedes v1--v5.}

\begin{document}

\twocolumn[{%
\centering
{\fontsize{15}{18}\selectfont\bfseries A License-Clean Graph Neural Network for Fast
Parasitic Extraction in PCB-Embedded Planar Magnetics\par}
\vspace{7pt}
{\fontsize{10}{12}\selectfont Duong Viet Hoang$^{1,*}$, \quad Lun-Min Shih$^{1}$\par}
\vspace{2pt}
{\fontsize{9}{11}\selectfont $^{1}$Department of Computer Science, Da-Yeh University, Taiwan\par}
{\fontsize{9}{11}\selectfont Hoangduong4316@icloud.com, \quad Lmshih@mail.dyu.edu.tw \quad ($^{*}$Corresponding author)\par}
\vspace{2pt}
{\fontsize{8}{9.5}\selectfont\itshape Version 3 (post-review): adds the gradient-boosted-tree baseline (Table~II) answering "why a graph?"; notation unified on the winding mutual $M$. Supersedes v1--v2.\par}
\vspace{7pt}
}]

\noindent\textbf{Abstract}---High-density power modules embed their planar magnetics
in the PCB, where extracting the parasitic capacitance and inductances from a
layout is a design bottleneck. Commercial 3-D extractors (e.g.\ ANSYS Q3D) are
accurate but closed and costly. We present a license-clean surrogate: a
geometry-aware message-passing graph neural network (275\,k parameters, pure
PyTorch, no GPL or torch-geometric dependency) that predicts the four lumped
quantities of 8-layer planar-LLC layouts: the inter-winding capacitance $C_{ps}$
--- the parasitic proper --- plus the air-core self/leakage inductances $L_p,L_s$
and mutual $M$. We
first diagnose the analytical labels as $52$-$59\%$ off \textsc{FastHenry}
(inductance) and $267\%$ off a 3-D electrostatic FEM ($C_{ps}$; its overlap is
position-blind). Regenerating every label with its 3-D solver and retraining
brings the GNN to $2.6$-$3.8\%$ on $L_p/L_s/M$ and $3.2\%$ on $C_{ps}$
($R^2\!\ge\!0.95$), so \emph{all four} targets are field-grade
against 3-D field solvers; an independent Neumann-integral reference puts the
GNN's error inside the inter-solver gap, and the core magnetizing term
($37$-$73\times$ the winding mutual) is validated separately against a real
commercial core's datasheet $A_L$ ($\sim$5\%). The value is a
\emph{speed-accuracy} one: the GNN returns this solver-grade vector in
$1.17$\,ms/design, $\sim$\textbf{4{,}300$\times$ faster} than the solvers it
matches, screening $1000$ interleavings in $\sim$1\,s rather than $\sim$1.4\,h,
where the fast analytical model is the wrong number and the accurate solver is
too costly per candidate. A pairwise ranking head orders unseen layout variants
at family-disjoint Spearman $\rho=0.93$; on the \emph{lateral-registration} lever
($x$-$y$ overlap) the analytical model is position-blind and cannot rank at all,
while the GNN tracks the 3-D solver at $\rho=0.95$ and its layout pick is
$\sim$$340\times$ lower-regret. We report two honest negatives ($E(n)$-equivariance
gives no gain; a dense graph is slower than analytical, though a sparse $k$-NN
model overtakes it at $N\!\approx\!80$). All accuracy is solver-grade, not
hardware-anchored; a measured-board campaign is the stated next step.

\vspace{3pt}
\noindent\textbf{Keywords}---PCB parasitic extraction; graph neural network; message
passing; PEEC; planar magnetics; power-electronics packaging; EMI.

\section{Introduction}
Planar magnetics realised as multilayer PCB foils are standard in GaN/SiC power
stages for their low profile, repeatable manufacturing, and thermal
spreading~\cite{ouyang,hurley}. Above roughly 1\,MHz the inter-winding
capacitance $C_{ps}$ and the self/leakage inductances dominate EMI and
switching-node ringing~\cite{ott,paul}, and the hand heuristics historically used
for spacing and interleaving no longer predict them~\cite{erickson}. Accurate
extraction requires 3-D field solvers; the de-facto standard, ANSYS Q3D, is
accurate but closed and carries a substantial per-seat licence cost, which
excludes it from the inner loop of an open design tool. A fast, license-clean
surrogate returning the parasitic matrix in milliseconds would let a layout
optimiser explore hundreds of interleavings per design.

We report a proof-of-concept of such a surrogate (Fig.~\ref{fig:pipeline}). Our
contributions are: (i) a geometry-aware MPNN predicting the four-target
graph-level parasitic vector with multi-seed accuracy ($R^2>0.999$), with a
non-trivial baseline and ablations isolating which design choices carry the
accuracy; (ii) a validation of \emph{all four} targets against 3-D field
solvers---after diagnosing the analytical labels as $52$-$59\%$ off
\textsc{FastHenry} and $267\%$ off a 3-D electrostatic FEM, we regenerate every
label with its solver and retrain to $2.6$-$3.8\%$ and $3.2\%$
($R^2\!\ge\!0.95$); the same solver family reproduces the datasheet $A_L$ of a
\emph{real commercial core} to $\sim$5\%, and the surrogate extends to a
frequency-resolved $R_{ac}/R_{dc}(f)$ at $0.37\%$; (iii) a diagnosis-and-fix for
ranking and a map of \emph{where} the learned ranker beats the analytical
baseline---a pairwise objective lifts $\rho=0.27$ to $\rho=0.93$, and on lateral
registration the position-blind analytical model cannot rank at all while the GNN
reaches $\rho=0.95$---plus an honest negative, an $E(n)$-equivariant architecture
does \emph{not} beat the invariant MPNN.\footnote{An earlier version of this work
mislabelled a linear regressor as a ``GNN'' and quoted an unreproducible
$183\times$ speedup; both are corrected here, and all numbers trace to released
run artifacts (see Reproducibility).}

\section{Related Work}
Message passing~\cite{gilmer} unifies a family of graph networks whose expressive
power follows the Weisfeiler-Lehman hierarchy~\cite{kipf,gin}; for physical
systems graph networks learn simulation dynamics directly~\cite{sanchez}, and
$E(n)$-equivariant networks~\cite{egnn} build rigid-motion symmetry into the
layers. In EDA~\cite{huang_survey}, GNNs address parasitic prediction
(ParaGraph~\cite{paragraph}), distributed-circuit EM behaviour~\cite{circuitgnn},
placement~\cite{mirhoseini}, and net-length estimation~\cite{net2}; these target
IC or RF \emph{netlists} and typically build on framework stacks (DGL/PyG). Our
setting differs in three ways: PCB-embedded \emph{planar magnetics} (foil traces,
coupled $R/L/C$ with turn-to-turn mutual inductance dominating); a deliberately
license-clean, dependency-light implementation (pure PyTorch~\cite{pytorch}, no
GPL, no torch-geometric); and supervision from a physically cross-validated
reference rather than commercial-tool labels. The
physics baseline is Ruehli's PEEC~\cite{ruehli} with Grover~\cite{grover} partial
inductances, Dowell~\cite{dowell} eddy/proximity resistance, and
Roshen~\cite{roshen} gap fringing, as compiled for magnetics
design~\cite{hurley,mclyman}. For field-level ground truth we use
scikit-fem~\cite{skfem} with Gmsh~\cite{gmsh}, cross-checked against surface-BEM
capacitance extraction~\cite{fastcap}. Operator-learning
surrogates~\cite{fno,deeponet} are the route to learning the full field rather
than lumped scalars.

\begin{figure*}[t]
\centering
\includegraphics[width=0.66\textwidth]{fig1_pipeline.pdf}
\caption{Pipeline. A planar PCB layout is encoded as a graph (nodes = trace
segments with invariant geometric features; edges = inductive/capacitive
couplings, dense all-pairs or sparse $k$-NN); a message-passing GNN regresses
$[C_{ps},L_p,L_s,M]$, supervised by 3-D field solvers.}
\label{fig:pipeline}
\end{figure*}

\section{Method}
\subsection{Graph construction, model, and setup}
Each rectangular foil trace is a node carrying rotation/translation-\emph{invariant}
features (width/length/thickness, $\log\sigma$, $\varepsilon_r$, normalized layer
index); coordinates enter only through relative edge vectors
$e_{ij}=(\lVert\Delta_{ij}\rVert,\mathrm{overlap},\Delta_{ij},\mathrm{type})$, so
the representation is invariant under rigid motion of the board. The \emph{dense}
graph connects all pairs ($O(N^2)$ edges), matching the reference's pair sum; the
\emph{sparse} graph keeps each node's $k$ nearest neighbours in 3-D ($O(Nk)$
edges), since parasitic coupling decays with distance. A forward pass costs
$O(|E|d^2)$ for hidden width $d$, so the dense model is $\Theta(N^2)$ and the
sparse model $\Theta(Nk)$---the asymptotic gap the scaling study
(Sec.~\ref{sec:scaling}) measures.

The network applies four edge-conditioned message-passing
layers~\cite{gilmer} (hidden width 96, SiLU, residual, LayerNorm) followed by a
mean\,$\oplus$\,max\,$\oplus$\,log-sum readout and a four-target head
(275{,}339 parameters); the log-sum channel matches the \emph{extensive} targets,
which scale with conductor count. Variable-size batching uses concatenated graphs
and segment aggregation, so the environment needs only PyTorch~\cite{pytorch}
(BSD). Training uses AdamW~\cite{adamw} ($10^{-5}$ decay), lr $2\times10^{-3}$
cosine, SmoothL1 on $\log$-standardized targets, 200 epochs, batch 32, seeds
$\{42,43,44\}$, 80/20 split with features standardized on the train split only.
The primary corpus is 2{,}000 synthetic 8-layer planar-LLC layouts (seed 42),
4-21 primary and 4-21 secondary foil traces interleaved across layers, averaging
25.2 nodes and 663 all-pairs edges per graph; labels are taken at the
low-frequency limit.

\subsection{Reference labels and their diagnosis}
The training target per layout is four \emph{extensive} scalars accumulated over
\emph{every} trace pair, an $O(N^2)$ operation, as a PEEC extractor builds its
partial-element matrix~\cite{ruehli}: $C_{ps}$ over all primary-secondary pairs,
$L_p/L_s$ as self plus intra-net mutual, and $M$ as the primary-secondary
mutual sum. \emph{What the solver returns, and what we predict.} \textsc{FastHenry}
does not report $L_p$ directly: with one port per trace it returns the
$N\!\times\!N$ \emph{partial}-inductance matrix, and the net value follows by the
series sum $L_p=\sum_i L_{ii}+2\sum_{i<j}L_{ij}$. The surrogate predicts that same
composition, so the comparison is like-for-like (the Neumann reference of
Sec.~\ref{sec:xsolver} uses the identical convention, different numerics). Scope,
then: $C_{ps}$ is the \textbf{parasitic} proper --- inter-winding capacitance,
sets the common-mode current, no functional role --- while the inductances are
\textbf{air-core} ($\mu_0$, no magnetic material), so $L_p$ is the winding's
free-space self/leakage inductance, \emph{not} the operating primary inductance,
which the ferrite magnetizing term dominates by $37$-$73\times$. We accordingly
write the coupling as the mutual $M$, not $L_m$. Because
these closed-form (``silver'') labels are analytical, we
cross-validate against independent field solves. Summing an electrostatic FEM
coupling over all primary-secondary pairs gives an independent total-$C_{ps}$
reference: over 12 layouts the analytical label sits $23.8\%$ from FEM and the
\emph{trained GNN} prediction $24.2\%$ from the same solve. The same test on the
dominant inductive output, against a multi-filament partial-inductance reference
(\textsc{FastHenry}-style~\cite{fasthenry}: each cross-section tiled into
$5\!\times\!2$ filaments, the pair mutual being the Neumann-integral average,
independent of Grover's closed form), gives $62.6\%$ for the analytical label and
$62.9\%$ for the GNN. This isolates the error source: the network is \emph{not}
the cause---it reproduces its teacher to $<\!1\%$---but the analytical Grover
mutual is poor against a partial-element reference. We act on this in
Sec.~\ref{sec:fieldgrade}: regenerating the labels is a label-quality fix, not an
architecture change.

\section{Results}
\subsection{Accuracy, baselines, and ablations}
Table~\ref{tab:acc} reports held-out accuracy of the dense GNN over three seeds:
all targets reach $R^2>0.999$ with small seed variance. Table~\ref{tab:abl}
answers the objection that a trivial regressor might match it. A pooled-feature
MLP (same readout, \emph{no} message passing) matches the GNN on $C_{ps}$ but
collapses on the inductances ($R^2\!\approx\!0.88$ vs.\ $0.9996$; RMSE $662$ vs.\
$37$\,nH on $L_p$). Zeroing the coordinates and edge relative-geometry features
roughly doubles the inductance RMSE, confirming that channel carries the inductive
information; a 2-layer model matches the 4-layer one, so accuracy comes from
message passing and geometry, not depth.

\emph{Why not gradient-boosted trees?} Table~\ref{tab:trees} runs XGBoost,
LightGBM, CatBoost, and RandomForest on the \emph{same} pooled 27-D vector,
labels, and a byte-identical 332/67 split. On the scalar capacitance the two are
level and the metrics split: the GNN wins on $R^2$ ($0.953$ vs.\ $0.942$) while
LightGBM wins on median relative error ($2.9\%$ vs.\ $3.2\%$). The disagreement is
informative rather than contradictory---$R^2$ is a squared-error measure dominated
by the largest-$C_{ps}$ layouts, the median is scale-free and outlier-robust---so
LightGBM is slightly better on a \emph{typical} layout while the GNN controls the
tail better. (RandomForest is the extreme case of the same effect: a $12.9\%$
median but $R^2=0.46$.) Either way that target alone does not justify a graph.
Trees lose on
everything relational: inductances $R^2=0.877$-$0.920$ vs.\ $0.994$-$0.995$;
ranking $\rho=0.088$ vs.\ $0.93$; and on size extrapolation their $M$ collapses to
$0.11$-$0.32$ where a GNN retrained on that split holds $0.943$. A fixed-dimension
tabular model also cannot be \emph{applied} to a different trace count without
re-engineering its features, whereas the GNN's shared message function transfers
by construction.

Table~\ref{tab:sweep} sweeps graph sparsity (reduced 150-epoch budget). Across
$k\in\{4,8,16,32\}$ every target stays at $R^2\ge0.979$ ($\ge0.99$ on $C_{ps}$):
accuracy is insensitive to sparsity because local coupling dominates, and even
$k=4$ ($\approx$100 edges/graph) suffices, which makes the $O(Nk)$ sparse model a
near-free speed-up rather than an accuracy trade (the mild non-monotonicity at
$k=8$ is within run-to-run variance at this single-seed budget). A capacity sweep
likewise shows a 32\,k-parameter model already at $R^2\approx0.999$, so the task
is not capacity-limited and the width-96 default is comfortably sufficient.

\emph{Size generalization.} Training on the smaller graphs ($\le32$ nodes) and
testing on the largest ($32$-$42$) gives $R^2 = 0.9988/0.9946/0.9957/0.9958$, a
mild degradation from the random-split $0.9993$-$0.9998$. The field-grade corpus
below extends the \emph{measured} range to 70 nodes ($R^2\ge0.988$); beyond that,
to the $N\!=\!1280$ of Sec.~\ref{sec:scaling}, accuracy is \emph{extrapolated}.

\begin{table}[t]
\caption{Dense-GNN accuracy, 3 seeds (held-out 400/2000, mean$\pm$std), and
baselines/ablations ($R^2$, seed 42, same reference).}
\label{tab:acc}
\centering
\small
\begin{tabular}{lrr}
\toprule
Target & RMSE & $R^2$\\
\midrule
$C_{ps}$ (pF)   & $23.1\pm0.9$  & $0.9998\,(\sigma\!<\!10^{-4})$\\
$L_{p}$ (nH)    & $37.5\pm2.7$  & $0.9996\,(\sigma\!<\!10^{-4})$\\
$L_{s}$ (nH)    & $39.3\pm5.6$  & $0.9995\pm0.0002$\\
$M$ (nH)    & $22.6\pm4.4$  & $0.9997\pm0.0001$\\
\bottomrule
\end{tabular}

\vspace{4pt}
\label{tab:abl}
\begin{tabular}{lrrrr}
\toprule
Model & $C_{ps}$ & $L_{p}$ & $L_{s}$ & $M$\\
\midrule
Dense GNN (4-layer)      & 0.9998 & 0.9996 & 0.9993 & 0.9995\\
GNN, 2-layer             & 0.9998 & 0.9996 & 0.9995 & 0.9994\\
GNN, no rel-geometry     & 0.9998 & 0.9985 & 0.9984 & 0.9985\\
Sparse $k$-NN GNN        & 0.9986 & 0.9927 & 0.9932 & 0.9983\\
Pooled MLP (no MP)       & 0.9879 & 0.8883 & 0.8802 & 0.9877\\
\bottomrule
\end{tabular}
\end{table}

\begin{table}[t]
\caption{Tree baselines vs.\ the GNN on the identical 332/67 field-grade split:
in-distribution $R^2$ (median relative error \%\ in parentheses), family-disjoint
ranking $\rho$, and the mutual-inductance $R^2$ on the size-extrapolation split.}
\label{tab:trees}
\centering
\footnotesize
\setlength{\tabcolsep}{3.5pt}
\begin{tabular}{lrrrr@{\ \ }rr}
\toprule
& \multicolumn{4}{c}{in-distribution $R^2$ (median \%)} & rank & size-gen\\
\cmidrule(lr){2-5}
Model & $C_{ps}$ & $L_{p}$ & $L_{s}$ & $M$ & $\rho$ & $R^2(M)$\\
\midrule
RandomForest & .462\,(12.9) & .887\,(24.8) & .914\,(15.9) & .880\,(9.3) & --- & 0.112\\
XGBoost      & .927\,(3.8)  & .878\,(16.3) & .908\,(16.5) & .914\,(7.8) & --- & 0.231\\
LightGBM     & .942\,(\textbf{2.9})  & .893\,(23.9) & .890\,(20.2) & .918\,(7.1) & 0.088 & 0.247\\
CatBoost     & .892\,(4.7)  & .896\,(24.5) & .893\,(20.4) & .920\,(8.1) & --- & 0.324\\
\midrule
\textbf{GNN} & \textbf{.953}\,(3.2) & \textbf{.995}\,(\textbf{2.8}) &
\textbf{.994}\,(\textbf{3.8}) & \textbf{.994}\,(\textbf{2.6}) &
\textbf{0.93} & \textbf{0.943}\\
\bottomrule
\end{tabular}
\end{table}

\begin{table}[t]
\caption{$k$-NN sparsity sweep ($R^2$, seed 42, 150-epoch budget).}
\label{tab:sweep}
\centering
\small
\begin{tabular}{rrrrrr@{\quad}}
\toprule
$k$ & edges & $C_{ps}$ & $L_{p}$ & $L_{s}$ & $M$\\
\midrule
4  & 101 & 0.9980 & 0.9911 & 0.9902 & 0.9976\\
8  & 201 & 0.9947 & 0.9804 & 0.9794 & 0.9947\\
16 & 397 & 0.9958 & 0.9805 & 0.9813 & 0.9965\\
32 & 643 & 0.9972 & 0.9945 & 0.9912 & 0.9951\\
\bottomrule
\end{tabular}
\end{table}

\subsection{Field-grade validation against 3-D solvers and a real core}
\label{sec:fieldgrade}
The diagnosis above motivates \emph{regenerating} the labels. We build a second
corpus (1{,}500 layouts, mean 39 / max 70 nodes) with $M$ labelled by the
multi-filament partial-element method, retrain, and validate the trained GNN's
total $M$ \emph{directly against \textsc{FastHenry}}~\cite{fasthenry}, the
established open-source 3-D partial-element inductance solver: our filament label
agrees with the full 3-D solve to \textbf{1.6\%} (median, 24 layouts) whereas the
Grover label is 56\% off, and the trained GNN matches the 3-D solver to
\textbf{2.2\%} (2.3\% mean)---field-grade against a 3-D field solver, not merely
against its own analytical label. An audit found the \emph{self}-inductance labels
equally poor ($L_p$ 51.5\%, $L_s$ 53.6\%, $M$ 58.6\%, median over 40 layouts),
so we relabel \emph{all three} with \textsc{FastHenry} and retrain: the GNN then
matches it to \textbf{2.8\,/\,3.8\,/\,2.6\%}.

\textsc{FastHenry} extracts inductance only, so $C_{ps}$ needs a capacitance
solver. Surface-BEM solvers (FastCap/FasterCap) did not converge on the
$\sim$0.1\,mm inter-layer gaps (uniform-panel FastCap swung 6-40$\times$ with
panel density; adaptive FasterCap cost $\sim$3\,min/solve), so we use a
\emph{volume} FEM solving $\nabla\!\cdot\!(\varepsilon\nabla\phi)=0$ on a gmsh
tetrahedral mesh of the air box with the real conductor boxes embedded (Dirichlet
$\phi=1/0$\,V on the primary/secondary windings), giving $C_{ps}=2W/V^2$. It is
robust (a conforming volume mesh has no BEM panel-orientation pathology), fast
($\sim$2\,s/solve), and \emph{geometry-correct}: on a stacked overlap it returns
92.9\,pF, on a 60\,mm-offset disjoint pair 0.38\,pF (a 244$\times$ spread),
whereas the position-blind analytical label gives the \emph{same} 44.3\,pF for
both. Relabelling and retraining, the GNN matches the FEM to \textbf{3.2\%}
($R^2=0.95$) against the analytical label's 267\%. All four targets are now
field-grade (Table~\ref{tab:fieldgrade}).

\subsection{Independent cross-solver check: is it physics or one solver's numerics?}
\label{sec:xsolver}
Training labels and held-out validation both come from the same solver family, so
a legitimate objection is that the GNN may have learned \textsc{FastHenry}'s
numerics rather than the physics. We test this against a \emph{third} reference
built independently of both: a direct Neumann double-integral partial inductance
with $5\!\times\!2$ filament subdivision, sharing no code path with the
\textsc{FastHenry} binary. Over 70 held-out layouts (346 total, seed 42) the GNN
matches this independent reference at $R^2=0.985/0.988/0.993$ and median
$6.6/5.0/2.0\%$ for $L_p/L_s/M$; \textsc{FastHenry} itself matches it at
$R^2=0.994/0.996/0.998$ and $5.0/4.8/2.1\%$. The GNN's error is therefore at or
inside the \emph{inter-solver} gap for $M$ and $L_s$: the surrogate sits within
the spread that two independent formulations of the same physics already show
between themselves, which is the appropriate resolution limit and not a
solver-specific artifact.

\emph{Solver check on a real commercial part.} On the EER40 GP95 ferrite from a
published PFC-choke worked example, the same magnetostatic FEM reproduces the
datasheet $A_L$ to $\sim$5\% at production gaps (5.7/4.1/5.4\% at
0.05/0.1/0.2\,mm). This validates the \emph{solver} against a measured datasheet,
not the GNN, which the EER40 case does not feed. Placing our layouts on that core
is what puts the magnetizing inductance $37$-$73\times$ above the air-core $M$.

\emph{Frequency-dependent extraction.} AC winding resistance rises with frequency
through skin and proximity effects, so we extend the surrogate from scalars to the
primary $R_{ac}/R_{dc}(f)$ curve over 10\,kHz-100\,MHz (17 points), ground-truthed
by a \textsc{FastHenry} sweep with filament subdivision (mean curve
$1.0\!\rightarrow\!6.0\times$). Trained on 136 small boards, the same GNN with a
17-dimensional head predicts held-out curves to \textbf{0.37\% median} error
(curve $R^2=0.995$), closing the ``low-frequency only'' gap.

\emph{Does equivariance help? It does not.} An $E(n)$-equivariant network
(EGNN~\cite{egnn}, pure PyTorch) is \emph{marginally worse} on all four targets
($R^2$ eq/MPNN: 0.9993/0.9997, 0.9986/0.9992, 0.9983/0.9995, 0.9996/0.9998), with
no advantage in a sample-efficiency sweep (200/800 graphs); its only benefit is
near rotation-invariance (0.33\% median drift under a random 3-D rotation).
Because the targets are aggregated invariant scalars, the MPNN's
relative-geometry edge features already encode the needed symmetry.

\begin{table}[t]
\caption{All four parasitics field-grade: the trained GNN vs.\ 3-D field solvers
(\textsc{FastHenry} for $L$, an electrostatic FEM for $C_{ps}$), and the
analytical labels they replace.}
\label{tab:fieldgrade}
\centering
\small
\begin{tabular}{lr}
\toprule
Analytical Grover $L_p/L_s/M$ vs.\ \textsc{FastHenry} 3-D        & 52/54/59\,\%\\
Analytical $C_{ps}$ vs.\ FEM-3D (position-blind overlap)          & 267\,\%\\
\midrule
\textbf{GNN $L_p/L_s/M$ vs.\ \textsc{FastHenry} 3-D}            & \textbf{2.8/3.8/2.6\,\%}\\
\textbf{GNN $C_{ps}$ vs.\ 3-D electrostatic FEM} ($R^2{=}0.95$)   & \textbf{3.2\,\%}\\
Real core (EER40 GP95) FEM vs.\ datasheet $A_L$, prod.\ gaps       & $\sim$5\,\%\\
Large-$N$ measured, test 51-70 nodes, min ($R^2$)                 & $\ge0.988$\\
\bottomrule
\end{tabular}
\end{table}

\subsection{Speed at solver-grade accuracy}
\label{sec:scaling}
The design-loop value is a \emph{speed-accuracy} one (Fig.~\ref{fig:pareto}).
There are three ways to get the parasitic vector for a candidate stack-up: the
analytical PEEC (0.78\,ms/design, but $52$-$267\%$ off the 3-D solvers); the 3-D
solvers \textsc{FastHenry}+FEM ($\sim$5\,s/design, the accurate reference); or the
GNN, which returns the same four parasitics at \emph{field-grade} accuracy in
\textbf{1.17\,ms/design}, $\sim$\textbf{4{,}300$\times$ faster than the solvers it
matches}. It is thus the only option that is \emph{both} fast and 3-D-accurate:
the fast analytical model is the wrong number, and the right number is three
orders of magnitude too slow for an inner loop. Screening $1000$ interleavings
costs $\sim$1.2\,s with the GNN versus $\sim$1.4\,h with the solvers, after which
only the top few candidates need a confirming 3-D solve.

\emph{Versus the fast baseline.} Even against the (fast but inaccurate) analytical
PEEC the GNN is competitive at scale: a dense all-pairs graph is $O(N^2)$ (1.17
vs.\ 0.78\,ms at the 25-node mean, the GNN slower), but a sparse $k$-NN graph
($k=8$, retrained, $R^2\ge0.992$) is $\approx O(N)$ and overtakes the
unaccelerated PEEC at $N\!\approx\!80$ segments, $\sim80\times$ at $N=1280$
(extrapolated accuracy there); a multipole PEEC would itself approach $O(N)$. The
headline claim remains the $4{,}300\times$ over the \emph{accurate} solvers, which
no fast analytical model matches for accuracy.

\begin{figure}[t]
\centering
\includegraphics[width=0.84\columnwidth]{fig_pareto.pdf}
\caption{Speed-accuracy trade-off (both axes log). The analytical PEEC is fast but
$52$-$267\%$ off the 3-D solvers; the solvers are accurate but $\sim$5\,s/design;
the GNN occupies the fast-and-accurate corner (3.2\% median at 1.17\,ms).}
\label{fig:pareto}
\end{figure}

\begin{figure}[t]
\centering
\includegraphics[width=0.84\columnwidth]{fig_lateral.pdf}
\caption{Why the GNN's ranking beats the analytical model on lateral registration.
(a) Shifting the secondary trace laterally shrinks the true P/S overlap. (b) The
analytical $C_{ps}$ uses $\min(L)\min(w)$ and is position-blind, so it stays
constant ($550$\,pF) and cannot rank; the 3-D FEM (and the GNN tracking it at
$\rho=0.95$) follows the real overlap and drops $67\%$.}
\label{fig:lateral}
\end{figure}

\section{Discussion and Limitations}
\label{sec:limits}
The accuracy is reproducible across seeds, the graph structure is shown to be
necessary, all four targets are validated against 3-D field solvers, the physics
is checked against a real commercial core, and the fast sparse model is itself
accuracy-validated. Four limitations remain. (i) Validation is solver-grade, not
hardware-grade: ``field-grade'' means matching a solver, not a measured board. The
cross-solver check (Sec.~\ref{sec:xsolver}) rules out single-solver overfitting but
not a shared modelling assumption, the scope is the winding terms rather than the
core magnetizing inductance, and the $C_{ps}$ FEM uses a
uniform effective permittivity. Full-wave/Q3D and silicon anchoring remain open.
(ii) The geometry is synthetic (random interleaved rectangular foils; no vias,
return planes, or board-level parasitics). (iii) An $E(n)$-equivariant
architecture did \emph{not} improve accuracy or sample efficiency; the gains here
come from better labels and a frequency target, not architectural symmetry.

(iv) \textbf{Ranking: a ranking objective, and the lever where it matters.}
\emph{First}, ranking needs a ranking objective: a pointwise regressor orders
within-family variants at only $\rho=0.27$ (family-disjoint), because at global
$R^2=0.95$ its error floor ($\sqrt{1-R^2}\cdot\sigma$) rivals the small
within-family $C_{ps}$ spread, so a globally-excellent regressor is not a
fine-grained ranker (an earlier $\rho=0.51$-$0.74$ was a leakage artifact). A
\emph{pairwise} margin-ranking objective on the FEM-3D labels, in the
learning-to-rank tradition~\cite{ranknet,lambdamart}, fixes this ($\rho=0.93$).
\emph{Second}, \emph{where} it beats the free analytical baseline: an optimiser
turns two unequal levers. For \textbf{interleaving} (P/S layer/$z$ order) the
analytical PEEC is already competitive ($\rho=0.78$ vs.\ the FEM; GNN $0.85$),
because the $z$-gap is in its parallel-plate formula. For \textbf{lateral
registration} (in-plane $x$-$y$ overlap, Fig.~\ref{fig:lateral}) the PEEC is
\emph{blind}: its overlap $\min(L)\min(w)$ is position-independent, so it returns
a \emph{constant} $C_{ps}$ across every test family and cannot rank at all, while
the FEM $C_{ps}$ varies by $\sim$$100\%$ and the GNN tracks it at $\rho=0.95$.
\emph{This changes the design choice}: picking the lowest-$C_{ps}$ layout per
family, the PEEC-screened pick carries $1176\%$ regret vs.\ the FEM optimum---no
better than random ($1126\%$), the screener being blind---whereas the
GNN-screened pick has $3.5\%$, a $\sim$$340\times$ better decision. On
interleaving both pick well ($5.4$ vs.\ $5.8\%$). The residual cap is that these
decisions are against the FEM, not a measured board.

\emph{Integration.} In the inner loop of a layout optimiser the regression head
returns absolute $[C_{ps},L_p,L_s,M]$ in $\sim$1\,ms and the pairwise head orders
unseen designs at $\rho=0.93$, so hundreds of candidates are screened before any
field solve; the $\sim$1\,MB pure-PyTorch checkpoint ships inside the open tool
and runs on a laptop CPU.

\section{Conclusion}
A license-clean geometry-aware GNN predicts the planar-PCB $C_{ps}$ and air-core
inductance vector at $R^2>0.999$ over three seeds, with accuracy shown to require message
passing and geometry. After diagnosing a $62.6\%$ Grover-label error and
regenerating every label with its 3-D solver, all four targets are field-grade
($2.6$-$3.8\%$ on $L_p/L_s/M$, $3.2\%$ on $C_{ps}$), an independent Neumann
reference places the GNN inside the inter-solver gap, and the separately-validated
core term is quantified at $37$-$73\times$ the winding mutual. We extend the
surrogate to $R_{ac}/R_{dc}(f)$ at $0.37\%$ and fix the ranking gap ($\rho=0.27
\rightarrow 0.93$, decisive on the lateral lever where the analytical model is
blind), and report two honest negatives ($E(n)$-equivariance gives no gain; the
dense graph is slower than analytical, the sparse $k$-NN variant overtaking it at
$N\!\approx\!80$ at $R^2\ge0.992$). Next: silicon/Q3D anchoring, vias and return
planes, and an operator-learning surrogate~\cite{fno,deeponet} for the full field.

\section*{Reproducibility}
The release includes the dataset generator (fixed seed), per-run logs for every
experiment reported here, the trained checkpoints, and the FEM and
\textsc{FastHenry} reference solvers (built from source, MIT-licensed). The code
(pure PyTorch/NumPy/SciPy/scikit-fem) is BSD-3; every number regenerates from
these artifacts.

{\scriptsize
\setlength{\itemsep}{-3pt}\setlength{\parskip}{-1pt}
\begin{thebibliography}{00}
\bibitem{ouyang} Z. Ouyang and M. A. E. Andersen, ``Overview of planar magnetic
technology---Fundamental properties,'' \emph{IEEE Trans.\ Power Electron.}, vol.\
29, no.\ 9, pp.\ 4888-4900, 2014.
\bibitem{hurley} W. G. Hurley and W. H. W\"olfle, \emph{Transformers and Inductors
for Power Electronics}. Wiley, 2013.
\bibitem{mclyman} C. W. T. McLyman, \emph{Transformer and Inductor Design
Handbook}, 4th ed. CRC Press, 2011.
\bibitem{erickson} R. W. Erickson and D. Maksimovi\'c, \emph{Fundamentals of Power
Electronics}, 3rd ed. Springer, 2020.
\bibitem{dowell} P. L. Dowell, ``Effects of eddy currents in transformer
windings,'' \emph{Proc.\ IEE}, vol.\ 113, no.\ 8, pp.\ 1387-1394, 1966.
\bibitem{roshen} W. A. Roshen, ``Fringing field formulas and winding loss due to
an air gap,'' \emph{IEEE Trans.\ Magn.}, vol.\ 43, no.\ 8, pp.\ 3387-3394, 2007.
\bibitem{ruehli} A. E. Ruehli, ``Equivalent circuit models for three-dimensional
multiconductor systems,'' \emph{IEEE Trans.\ Microw.\ Theory Tech.}, vol.\ 22, no.\
3, pp.\ 216-221, 1974.
\bibitem{fasthenry} M. Kamon, M. J. Tsuk, and J. K. White, ``FASTHENRY: A
multipole-accelerated 3-D inductance extraction program,'' \emph{IEEE Trans.\
Microw.\ Theory Tech.}, vol.\ 42, no.\ 9, pp.\ 1750-1758, 1994.
\bibitem{grover} F. W. Grover, \emph{Inductance Calculations: Working Formulas and Tables}. Van Nostrand, 1946 (repr.\ Dover, 1962).
\bibitem{ott} H. W. Ott, \emph{Electromagnetic Compatibility Engineering}. Wiley,
2009.
\bibitem{paul} C. R. Paul, \emph{Introduction to Electromagnetic Compatibility},
2nd ed. Wiley, 2006.
\bibitem{skfem} T. Gustafsson and G. D. McBain, ``scikit-fem: A Python package for
finite element assembly,'' \emph{J.\ Open Source Softw.}, vol.\ 5, no.\ 52, 2369,
2020.
\bibitem{gmsh} C. Geuzaine and J.-F. Remacle, ``Gmsh: A 3-D finite element mesh
generator with built-in pre- and post-processing facilities,'' \emph{Int.\ J.\
Numer.\ Methods Eng.}, vol.\ 79, no.\ 11, pp.\ 1309-1331, 2009.
\bibitem{pytorch} A. Paszke \emph{et al.}, ``PyTorch: An imperative style,
high-performance deep learning library,'' in \emph{NeurIPS}, 2019.
\bibitem{gilmer} J. Gilmer, S. S. Schoenholz, P. F. Riley, O. Vinyals, and G. E.
Dahl, ``Neural message passing for quantum chemistry,'' in \emph{ICML}, 2017,
pp.\ 1263-1272.
\bibitem{kipf} T. N. Kipf and M. Welling, ``Semi-supervised classification with
graph convolutional networks,'' in \emph{ICLR}, 2017.
\bibitem{gin} K. Xu, W. Hu, J. Leskovec, and S. Jegelka, ``How powerful are graph
neural networks?'' in \emph{ICLR}, 2019.
\bibitem{egnn} V. G. Satorras, E. Hoogeboom, and M. Welling, ``E(n) equivariant
graph neural networks,'' in \emph{ICML}, 2021.
\bibitem{sanchez} A. Sanchez-Gonzalez \emph{et al.}, ``Learning to simulate complex
physics with graph networks,'' in \emph{ICML}, 2020, pp.\ 8459-8468.
\bibitem{huang_survey} G. Huang \emph{et al.}, ``Machine learning for electronic
design automation: A survey,'' \emph{ACM Trans.\ Des.\ Autom.\ Electron.\ Syst.},
vol.\ 26, no.\ 5, Art.\ 40, pp.\ 40:1-40:46, 2021.
\bibitem{paragraph} H. Ren \emph{et al.}, ``ParaGraph: Layout parasitics and
device parameter prediction using graph neural networks,'' in \emph{DAC}, 2020.
\bibitem{circuitgnn} G. Zhang, H. He, C.-Y. Hsu, and D. Katabi, ``Circuit-GNN: Graph
neural networks for distributed circuit design,'' in \emph{ICML}, 2019, pp.\
7364-7373.
\bibitem{mirhoseini} A. Mirhoseini \emph{et al.}, ``A graph placement methodology
for fast chip design,'' \emph{Nature}, vol.\ 594, pp.\ 207-212, 2021.
\bibitem{net2} Z. Xie \emph{et al.}, ``Net2: A graph attention network method
customized for pre-placement net length estimation,'' in \emph{ASP-DAC}, 2021.
\bibitem{fastcap} K. Nabors and J. White, ``FastCap: A multipole accelerated 3-D
capacitance extraction program,'' \emph{IEEE Trans.\ Comput.-Aided Design},
vol.\ 10, no.\ 11, pp.\ 1447-1459, 1991.
\bibitem{fno} Z. Li \emph{et al.}, ``Fourier neural operator for parametric
partial differential equations,'' in \emph{ICLR}, 2021.
\bibitem{deeponet} L. Lu, P. Jin, G. Pang, Z. Zhang, and G. E. Karniadakis,
``Learning nonlinear operators via DeepONet based on the universal approximation
theorem of operators,'' \emph{Nat.\ Mach.\ Intell.}, vol.\ 3, pp.\ 218-229, 2021.
\bibitem{ranknet} C. J. C. Burges \emph{et al.}, ``Learning to rank using gradient
descent,'' in \emph{ICML}, 2005.
\bibitem{lambdamart} C. J. C. Burges, ``From RankNet to LambdaRank to LambdaMART: An
overview,'' Microsoft Research, Tech.\ Rep.\ MSR-TR-2010-82, 2010.
\bibitem{adamw} I. Loshchilov and F. Hutter, ``Decoupled weight decay
regularization,'' in \emph{ICLR}, 2019.
\end{thebibliography}}

\end{document}
